% Part: first-order-logic 
% Chapter: axiomatic-deduction 
% Section: deduction-theorem-quantifiers

\documentclass[../../../include/open-logic-section]{subfiles}

\begin{document}

\olfileid[hi]{fol}{axd}{ddq}

\olsection{परिमाणकों सहित निगमन प्रमेय}

\begin{thm}[निगमन प्रमेय]
\ollabel{thm:deduction-thm-q} यदि $\Gamma \cup \{!A\} \Proves !B$,
तो $\Gamma \Proves !A \lif !B$।
\end{thm}

\begin{proof}
हम फिर $!B$ की, $\Gamma \cup \{!A\}$ से, !!{derivation} की लंबाई पर आगमन
करते हैं।

आगमन-आधार का प्रमाण~\olref[ded]{thm:deduction-thm} के प्रमाण वाले आधार के
समान है।

आगमन-चरण में फिर मान लें कि $!B$ की, $\Gamma \cup \{!A\}$ से,
!!{derivation} ऐसे चरण~$!B$ पर समाप्त होती है जिसे कोई अनुमान-नियम उचित
ठहराता है। यदि अनुमान-नियम पूर्वपक्ष-स्थापन है, तो
\olref[ded]{thm:deduction-thm} के प्रमाण की तरह आगे बढ़ते हैं। यदि
अनुमान-नियम \QR है, तो हम जानते हैं कि $!B \ident !C \lif \lforall[x][!D(x)]$ और
$!C \lif !D(a)$ रूप का !!a{formula},~!!{derivation}
में पहले आता है, जहाँ $a$,~$!C$, $!A$, या $\Gamma$ में नहीं आता। अतः हमारे
पास है:
\begin{align*}
  \Gamma \cup \{!A\} & \Proves !C \lif !D(a),\\
  \intertext{और आगमन-परिकल्पना लागू होती है; अर्थात् हमारे पास है}
    \Gamma & \Proves !A \lif (!C \lif !D(a)).\\
  \intertext{इससे}
  & \Proves (!A \lif (!C \lif !D(a))) \lif ((!A \land !C) \lif !D(a))\\
  \intertext{और पूर्वपक्ष-स्थापन द्वारा मिलता है}
  \Gamma & \Proves (!A \land !C) \lif !D(a).\\
  \intertext{क्योंकि अभिलाक्षणिक-चर शर्त अब भी लागू है, इसलिए इस !!{derivation} में \QR से उचित ठहराया गया चरण जोड़कर मिलता है}
    \Gamma & \Proves (!A \land !C) \lif \lforall[x][!D(x)].\\
    \intertext{हमारे पास यह भी है}
    & \Proves ((!A \land !C) \lif \lforall[x][!D(x)]) \lif (!A \lif (!C \lif \lforall[x][!D(x)]),\\
    \intertext{अतः पूर्वपक्ष-स्थापन से}
    \Gamma & \Proves !A \lif (!C \lif \lforall[x][!D(x)]),
\end{align*}
अर्थात् $\Gamma \Proves !B$।

जिस प्रसंग में $!B$ को \QR नियम उचित ठहराता है, पर वह
$\lexists[x][!D(x)] \lif !C$ रूप का है, उसे अभ्यास के लिए छोड़ते हैं।
\end{proof}

\begin{prob}
  \olref[fol][axd][ddq]{thm:deduction-thm-q} का प्रमाण पूरा कीजिए।
\end{prob}

\end{document}
